% !TeX root = ../教材题库_全部条目_转写工作稿.tex
% Source: XJ-B1-C05-TH-P157-Y0550
% 第5章 / 5.1 任意角和弧度制 / 教材 p.157 / PDF 第164页
% OCR未校对：公式与文字均可能有误；未完成数学TeX转写；禁止直接发给学生
% 原题及相邻原解截图（核对用）：../题目原图/XJ-B1-C05-TH-P157-Y0550-01.webp
% 本文件只在转写完成后作为题干使用；OCR 条目尚不是可靠题干。
\begin{bankitem}{p.157 思考（含紧邻说明）}{第5章 / 5.1 任意角和弧度制 / 教材 p.157 / PDF 第164页}{XJ-B1-C05-TH-P157-Y0550}{OCR未校对：公式与文字均可能有误；未完成数学TeX转写；禁止直接发给学生}
\begin{ocrdraft}
[教材p.157 / PDF第164页]
\par
？思考如图5．1一2，对于平面直角坐标系内任意一条射线OB，以OB为终边
\par
的角仅是否唯一？如果不唯一，那么终边相同的角有什么关系？
\par
从角的形成过程可以看到，与某一个角的始边相同且终边相同的角有无数个，
\par
它们的大小相差360。的整数倍．例如，在图5，1一2中，405。角与45。角的终边重合，
\par
这两个角的大小之差为360。；而一240。角与120。角的终边重合，这两个角的大小之
\par
差为一360。．
\par
把所有与角a终边相同的角用集合表示出来，即
\par
0一十\ensuremath{\cdot}360。，k\ensuremath{\text{[OCR:EUR]}}Z\}，
\par
其中当差一0时，角就是角仅本身．
\end{ocrdraft}
\end{bankitem}
